\chapter{The \code{rosetta\_gen} tool}
\label{ch:rosettagen}

\code{rosetta\_gen} is the single binary that drives the whole pipeline. It
reads a manifest and can take you from ``a folder of C++ headers'' to ``compiled
bindings for every declared language'' in one command --- or one step at a time,
when you want to inspect what each stage produces.

The tool itself uses \textbf{no reflection}, only JSON parsing and text
templating, so it builds with an off-the-shelf C++17 compiler. Only the
\emph{generator} it emits --- and the \lang{rest} and \lang{json} targets ---
needs the clang-p2996 toolchain.

\section{Building it}

One-time, from the repository root (CMake $\geq$ 3.28; it fetches the
single-header nlohmann/json automatically):

\begin{lstlisting}[style=sh]
cmake -G Ninja -S tools/rosetta_gen -B tools/rosetta_gen/build
cmake --build tools/rosetta_gen/build
\end{lstlisting}

The binary always lands in \textbf{\code{<repo>/bin/rosetta\_gen}}, not in the
build tree. If you scaffolded your project with \code{tools/rosetta\_init.py},
its bootstrap \code{CMakeLists.txt} does this for you and the binary appears at
\code{extern/rosetta/bin/rosetta\_gen}.

\section{The four modes}

\begin{lstlisting}[style=plain]
rosetta_gen <manifest.json> [out_dir]        # plain: emit the generator project only
rosetta_gen --build <manifest.json> [opts]   # the whole pipeline in one command
rosetta_gen --clean <manifest.json> [opts]   # remove everything generated
rosetta_gen --init  [manifest.json] [src]    # write a starter manifest
\end{lstlisting}

\code{rosetta\_gen --help} prints the full built-in help; \code{--build --help}
and \code{--clean --help} list each mode's options.

\begin{lstlisting}[style=plain]
--init --> manifest.json --> --build ------------> bindings/python  node  wasm ...
 (start)    (you edit it)      |                          ^
                               | = the three manual steps:|
                               |   1. emit gen/           |
                               |   2. cmake gen/, run it -+
                               |   3. cmake / npm / emcmake each backend
--clean <-- removes gen/ (or generated/), the generator binary and bindings/
\end{lstlisting}

\section{\code{--init}: starting a project}
\label{sec:init}

Writes a starter manifest (default \code{./manifest.json}). It \textbf{never
overwrites} an existing file --- remove it first to regenerate.

\subsection*{A blank, commented manifest}

\begin{lstlisting}[style=sh]
bin/rosetta_gen --init
\end{lstlisting}

Emits a fully-commented example exercising every commonly-used field: a
\field{variables} declaration feeding the \field{cpp26\_*} overrides, a
multi-entry \field{user\_include}, \field{user\_sources},
\field{compile\_definitions}, \field{build\_type} / \field{optimization}, a
representative spread of targets, one class and one function. The idea is that
you delete what you do not need rather than hunting the documentation for what
exists. \field{doc\_comments} is in there too, spelled at its default
(\code{true}) so you can see the switch exists --- it is the one field you set
in order to get \emph{less}.

\subsection*{Pre-filled from a source scan}

\begin{lstlisting}[style=sh]
bin/rosetta_gen --init manifest.json ./src
\end{lstlisting}

An argument naming an existing \textbf{directory} is the source tree to scan;
the other, if any, is the manifest path. The scan walks for headers
(\code{.h/.hh/.hpp/.hxx}) and sources (\code{.cpp/.cxx/.cc}) and pre-fills:

\begin{itemize}
  \item \field{classes} and \field{functions}, from the declarations found, with
    shared \field{namespace} / \field{header\_dir} defaults factored out when
    every entry agrees;
  \item \field{user\_sources} --- every source file found, \textbf{except} any
    defining \code{main()} (it would clash with the binding modules' entry
    points; reported as a note);
  \item \field{rosetta\_include}, guessed by walking up from the
    \code{rosetta\_gen} binary; a \code{FIXME} comment otherwise.
\end{itemize}

\begin{warn}
The scan is a comment- and string-aware \emph{token scan}, not a real C++ parse.
Template classes, overloaded free functions and anonymous namespaces are
skipped, each with a printed note. An overloaded free function is skipped
because the scan has no signature to write, not because it cannot be bound ---
add it by hand with a \field{signature} (Section~\ref{sec:signature}).
\textbf{Review the result before building.}
\end{warn}

VCS directories, build trees and vendored code (\code{build/}, \code{extern/},
\code{third\_party/}, \code{node\_modules/}, \ldots) are not descended into.

\begin{note}
For a whole \emph{project} skeleton --- a bootstrap \code{CMakeLists.txt} that
fetches rosetta into \code{extern/}, a \code{.gitignore}, a \code{README.md} ---
rather than just the manifest, use \code{tools/rosetta\_init.py}.
\end{note}

\section{\code{--build}: the whole pipeline}

\begin{lstlisting}[style=sh]
bin/rosetta_gen --build manifest.json
\end{lstlisting}

runs, in order:

\begin{enumerate}
  \item \textbf{emit} the generator project into \code{<manifest dir>/gen/}.
    This is also where the bound headers are read for their doc comments and
    default arguments (chapter~\ref{ch:documentation}), so the driver carries
    them as data --- no toolchain involved, and \field{doc\_comments}
    \code{false} skips it;
  \item \textbf{build} it with CMake --- this step needs the clang-p2996
    toolchain, and reflection runs here --- and \textbf{run} the resulting
    generator, writing one self-contained project per target into
    \code{<manifest dir>/bindings/<lang>/};
  \item \textbf{compile every backend} the manifest declares, each with its own
    build shape.
\end{enumerate}

\begin{center}
\small
\begin{tabular}{@{}L{6.6cm}L{6.4cm}@{}}
\toprule
\textbf{Backend} & \textbf{Build} \\
\midrule
\lang{python}, \lang{nanobind}, \lang{rest}, \lang{json}, \lang{lua},
  \lang{imgui}, \ldots & plain \code{cmake} configure + build \\
\lang{node}   & \code{npm install} + \code{npm run build} \\
\lang{wasm}   & \code{emcmake cmake} + \code{cmake --build} \\
\lang{julia}  & \code{cmake} (locates JlCxx by running \code{julia}) \\
\lang{csharp} & \code{cmake}, then \code{dotnet build} as a second stage \\
\lang{java}   & \code{cmake}, then \code{mvn package} as a second stage \\
\lang{qt}, \lang{qml} & \code{cmake} with \code{-DQT\_DIR} \\
\lang{markdown}, \lang{html}, \lang{typescript}, \lang{openapi},
  \lang{paraview} & nothing to compile \\
\bottomrule
\end{tabular}
\end{center}

A backend whose toolchain is \textbf{missing} --- no \code{npm}, no
\code{emcc}, no \code{julia} --- is \emph{skipped with a note}, not an error.
Likewise a missing \code{dotnet} or \code{mvn} downgrades \lang{csharp} /
\lang{java} to ``OK (native only)''. A backend whose build \textbf{fails} is
reported, and \code{--build} exits non-zero \emph{after trying the rest}.

Every run ends with a per-backend summary:

\begin{lstlisting}[style=plain]
== summary
   python      OK
   node        OK
   wasm        SKIPPED (emcc not found -- activate emsdk)
   markdown            OK (nothing to compile)
\end{lstlisting}

\subsection{Options}

\begin{lstlisting}[style=plain]
--gen-dir DIR            generator project dir   (default: <manifest dir>/gen)
--bindings-dir DIR       generated bindings dir  (default: <manifest dir>/bindings)
--clang-p2996-root PATH  -DCLANG_P2996_ROOT for every CMake configure
--qt-dir PATH            -DQT_DIR for the qt / qml builds
--only a,b / --skip a,b  restrict the backend builds to / exclude these
--cmake-arg ARG          extra argument for every CMake configure (repeatable)
--jobs N, -jN            parallel build jobs (bare -j: one per core)
--fresh                  wipe the gen and bindings dirs first
--wheel                  also build a Python wheel for the python /
                         nanobind targets
--wheel-dir DIR          collect every wheel in DIR (implies --wheel)
\end{lstlisting}

\subsection{Iterating on one or two backends}
\label{sec:onlyskip}

While developing you rarely want the full matrix. \code{--only} and
\code{--skip} filter \textbf{which backends get compiled} --- generation still
emits them all, and the filters name the \field{lang} values:

\begin{lstlisting}[style=sh]
# just python and node, 8-way parallel
bin/rosetta_gen --build manifest.json --only python,node -j8

# everything except the slow wasm link
bin/rosetta_gen --build manifest.json --skip wasm
\end{lstlisting}

\subsection{A custom reflection toolchain}

\lang{rest} is the one target whose generated code still splices reflections, so
it re-runs reflection when \emph{it} compiles. Its CMake configure needs to know
where the clang-p2996 build lives if it is not at the manifest's default:

\begin{lstlisting}[style=sh]
bin/rosetta_gen --build manifest.json \
    --clang-p2996-root ~/devs/clang-p2996/build --fresh
\end{lstlisting}

\code{--fresh} wipes \code{gen/} and \code{bindings/} first --- the right move
after toolchain or manifest-layout changes, when stale CMake caches would
otherwise bite.

\subsection{Qt and QML}

\begin{lstlisting}[style=sh]
bin/rosetta_gen --build manifest.json --only qt,qml --qt-dir ~/Qt/6.8.3/macos
\end{lstlisting}

Equivalently, set \field{qt\_dir} in the manifest and drop the flag.

\subsection{Passing flags to every configure}

\code{--cmake-arg} is repeatable and applies to the generator's configure
\emph{and} every backend's, including the emcmake configure for wasm:

\begin{lstlisting}[style=sh]
bin/rosetta_gen --build manifest.json \
    --cmake-arg -DCMAKE_BUILD_TYPE=Debug \
    --cmake-arg -DCMAKE_EXPORT_COMPILE_COMMANDS=ON
\end{lstlisting}

\subsection{Python wheels}

By default \code{--build} compiles the extension module and stops --- it
produces an importable \code{.so} / \code{.pyd}, not a distributable artifact.
\code{--wheel} adds the packaging step for the two backends that emit a
\code{make\_wheel.py}:

\begin{lstlisting}[style=sh]
bin/rosetta_gen --build manifest.json --only python,nanobind --wheel
\end{lstlisting}

Wheels land in \code{<bindings>/<lang>/dist}, with the repaired
(redistributable) copies in \code{dist/repaired}. It is opt-in because
packaging is a release action, not a development one: it pip-installs
\code{build}, \code{scikit-build-core} and a repair tool into the environment,
and takes noticeably longer than the compile \code{--build} is otherwise
optimised for.

\code{--wheel-dir} collects every backend's wheels in one place instead --- the
shape you want when the next step is uploading them:

\begin{lstlisting}[style=sh]
bin/rosetta_gen --build manifest.json --wheel-dir ./dist
\end{lstlisting}

\begin{lstlisting}[style=plain]
dist/
  pygeom-0.2.0-cp312-cp312-macosx_15_0_arm64.whl
  nbgeom-0.2.0-cp312-abi3-macosx_15_0_arm64.whl
  repaired/          <- the redistributable copies, same names
\end{lstlisting}

Sharing one output directory is safe: each \code{make\_wheel.py} repairs only
the wheels \emph{it} just produced.

Both flags have manifest counterparts --- \field{wheel} and
\field{wheel\_dir} on the \lang{python} / \lang{nanobind} entry of
\field{targets} (Section~\ref{sec:wheels}) --- so a project that always
packages says so once. They are per-target, which is what lets a manifest ship
the \lang{nanobind} wheel and not the \lang{python} one. The flags still win,
and only ever toward packaging \emph{more}: a target spelling
\code{"wheel": false} cannot disarm a \code{--wheel} you typed, and
\code{--wheel-dir} overrides every target's own directory.

\begin{note}
The wheel build is \emph{independent} of the compile above it ---
scikit-build-core re-runs the same \code{CMakeLists.txt} in its own tree, so it
neither reuses nor disturbs \code{<lang>/build/}. The interpreter is probed on
\code{PATH}; if none is found, the backend is recorded
\code{OK (no wheel -- no python interpreter found)} rather than failing. A
\code{--wheel} that cannot fire --- no such target, or none surviving
\code{--only}/\code{--skip} --- says so instead of quietly doing nothing.
\end{note}

\section{Plain mode: one step at a time}

\begin{lstlisting}[style=sh]
bin/rosetta_gen manifest.json [out_dir]
\end{lstlisting}

emits \textbf{only} the generator project (default \code{out\_dir}:
\code{<manifest dir>/generated/}) and stops. Useful for inspecting
\code{bindings.h}, hacking on the emitted driver, or wiring the steps into your
own build system:

\begin{lstlisting}[style=sh]
# 1. manifest -> generator project
bin/rosetta_gen manifest.json gen

# 2. build the generator (clang-p2996) and run it
cmake -S gen -B gen/build && cmake --build gen/build
./generator bindings          # the binary is dropped NEXT TO the project dir

# 3. build any backend by hand
cmake -S bindings/python -B bindings/python/build
cmake --build bindings/python/build
\end{lstlisting}

Two footguns the tool warns about:

\begin{itemize}
  \item a fresh checkout (or a \code{--clean}) leaves \code{out\_dir/build}
    unconfigured, so \code{cmake --build} alone fails --- the \emph{configure}
    step is printed as a \code{next:} hint;
  \item if the emitted sources changed but a generator binary from a previous
    emit still sits next to \code{out\_dir}, running it without rebuilding
    \textbf{silently re-emits the old bindings}. The tool prints a rebuild note
    when it detects this.
\end{itemize}

Re-emitting is idempotent: unchanged files are not rewritten, so the next
\code{cmake --build} recompiles nothing.

\section{\code{--clean}: back to sources}

\begin{lstlisting}[style=sh]
bin/rosetta_gen --clean manifest.json [--gen-dir DIR] [--bindings-dir DIR]
\end{lstlisting}

Removes everything the pipeline generated for this manifest:

\begin{itemize}
  \item the generator project directory --- the explicit \code{--gen-dir}, or
    \textbf{both} defaults (\code{gen/} from \code{--build} and
    \code{generated/} from plain mode);
  \item the generator binary dropped beside it, including per-config
    subdirectories of multi-config generators;
  \item the per-backend folders under \code{bindings/} \textbf{that the manifest
    declares}. Anything else you parked there survives, and \code{bindings/}
    itself is only removed once empty.
\end{itemize}

Never the manifest --- and never a folder whose \code{CMakeLists.txt} lacks the
\code{Generated by rosetta} stamp the emitters write, so a hand-written folder
sharing a name survives:

\begin{lstlisting}[style=plain]
not removing gen (no rosetta_gen marker in its CMakeLists)
\end{lstlisting}

\section{Exit status}

\begin{center}
\small
\begin{tabular}{@{}lL{4.6cm}L{5.6cm}@{}}
\toprule
\textbf{Mode} & \textbf{0} & \textbf{non-zero} \\
\midrule
plain & project emitted & manifest missing / invalid, emit error \\
\code{--build} & every \emph{attempted} backend built (skips do not count) &
  generator build or run failed, or at least one backend build failed \\
\code{--clean} & clean ran (nothing to remove is fine) &
  manifest missing / invalid \\
\code{--init} & manifest written & target file already exists, bad arguments \\
\bottomrule
\end{tabular}
\end{center}
